%% LyX 2.4.4 created this file.  For more info, see https://www.lyx.org/.
%% Do not edit unless you really know what you are doing.
\documentclass[english]{article}
\usepackage[T1]{fontenc}
\usepackage[latin9]{inputenc}
\usepackage{enumitem}
\usepackage{amsmath}
\usepackage{amsthm}
\usepackage{cancel}
\usepackage{geometry}
\geometry{verbose,tmargin=1in,bmargin=1in,lmargin=1in,rmargin=1in}

\makeatletter

%%%%%%%%%%%%%%%%%%%%%%%%%%%%%% LyX specific LaTeX commands.
%% Because html converters don't know tabularnewline
\providecommand{\tabularnewline}{\\}

%%%%%%%%%%%%%%%%%%%%%%%%%%%%%% Textclass specific LaTeX commands.
\numberwithin{equation}{section}
\numberwithin{figure}{section}
\newlength{\lyxlabelwidth}      % auxiliary length 

%%%%%%%%%%%%%%%%%%%%%%%%%%%%%% User specified LaTeX commands.
\usepackage{hyperref}
\usepackage[usenames,dvipsnames,table]{xcolor} %adds additional color options
\hypersetup{colorlinks=true, urlcolor=black}

\usepackage{fancyhdr}
\usepackage{lastpage}
\pagestyle{fancy}
\usepackage{enumitem}
\usepackage{ifthen}
%sets math fonts to be more readable
\everymath=\expandafter{\the\everymath\displaystyle}
%\DeclareMathSizes{10}{10}{10}{10}
\usepackage{multicol}

\usepackage{tikz}
\usetikzlibrary{plotmarks}
\usepackage{pgfplots}
\pgfplotsset{compat=newest}

\usepackage{calc}
\usepackage[nomessages]{fp}

\usepackage{advdate}    % Advancing/saving dates
\usepackage{datetime}   % Dates formatting
\usepackage{datenumber} % Counters for dates

\usepackage{fancybox} %%ovalbox and fancy box settings
% Added by lyx2lyx
\setlength{\parskip}{\smallskipamount}
\setlength{\parindent}{0pt}

\makeatother

\usepackage{babel}
\begin{document}
\renewcommand{\author}{Dr.\,James Ashe}

\newcommand{\classNum}{107}

\newcommand{\classSec}{7 and 10}

\newcommand{\nameType}{\hspace{2cm}Name:}

\newcommand{\examType}{Quiz 3: 1.5}

\renewcommand{\date}{\formatdate{26}{9}{2014}}

%%First page's header%%%%%%%%%%%%%%%%%%%%%%
\fancypagestyle{firststyle} %%header settings for first pagr
{
	\fancyhead[L]{MTH \classNum{}(\classSec{}) \author{}\\
	\textbf{\examType{}} \date{}}
	\fancyhead[C]{\textbf{\nameType}}
	\setlength{\headsep}{14pt}
}
%%%%%%%%%%%%%%%%%%%%%%%%%%%%%%%%%%%%%%%%%%

%%Next page's header%%%%%%%%%%%%%%%%%%%%%%%
%\setlist{nolistsep}
\lhead{MTH \classNum{}(\classSec{})} 
\chead{\textbf{\nameType}} 
\cfoot{\thepage{}~of~\pageref{LastPage}} 
\setlength{\headsep}{5pt} 
%%%%%%%%%%%%%%%%%%%%%%%%%%%%%%%%%%%%%%%%%%%

%%Various settings; see the preamble for other settings%%%%%
%%\setenumerate[0]{leftmargin=12pt,itemindent=0pt,labelwidth=0pt}
\renewcommand{\labelenumi}{\textbf{\arabic{enumi}.}}
\renewcommand{\labelenumii}{\textbf{\alph{enumii}.}}
\thispagestyle{firststyle}
%%%%%%%%%%%%%%%%%%%%%%%%%%%%%%%%%%%%%%%%%%
\newcommand{\xmax}{10}
\newcommand{\xmin}{-10}
\newcommand{\xstep}{0} %must be xmin+n
\newcommand{\ymax}{10}
\newcommand{\ymin}{-10}
\newcommand{\ystep}{0} %must be ymin+n

\newcommand{\xspace}{\hspace*{40pt}}
\newcommand{\yspace}{\vspace*{.75cm}}
\newcommand{\rowColor}{gray!40}
\large
\begin{enumerate}[resume]
\item Use a truth table to determine whether the arguments are valid.

\begin{enumerate}
\item If the defendant is innocent, he does not go to jail. The defendant
goes to jail. Therefore, the defendant is guilty.

\begin{enumerate}
\item []\vspace{-25pt}%
\begin{tabular}{rll}
 &  & \tabularnewline
 &  & \tabularnewline
\hspace*{10pt} & \hspace*{2.5cm} & 1. \tabularnewline
\hspace*{10pt} & \hspace*{2.5cm} & 2.\hspace*{4.5cm}\tabularnewline
\cline{3-3}
 &  & \tabularnewline
\end{tabular}
\item [] \hspace{-1.25cm}%
\begin{tabular}{c|c|c|c|c|c|c}
\multicolumn{1}{c}{} & \multicolumn{1}{c}{} & \multicolumn{1}{c}{} & \multicolumn{1}{c}{} & \multicolumn{1}{c}{} & \multicolumn{1}{c}{} & \tabularnewline
\multicolumn{1}{c}{} & \multicolumn{1}{c}{} & \multicolumn{1}{c}{\ovalbox{$1$}} & \multicolumn{1}{c}{\ovalbox{$2$}} & \multicolumn{1}{c}{} & \multicolumn{1}{c}{\ovalbox{$c$}} & \ovalbox{Argument}\tabularnewline
 &  & \xspace & \xspace & \xspace & \xspace & \xspace\tabularnewline
\hline 
\hline 
\rowcolor{\rowColor}T & T &  &  &  &  & \tabularnewline
\hline 
T & F &  &  &  &  & \tabularnewline
\hline 
\rowcolor{\rowColor}F & T &  &  &  &  & \tabularnewline
\hline 
F & F &  &  &  &  & \tabularnewline
\end{tabular}\vfill
\end{enumerate}
\end{enumerate}
\begin{enumerate}[resume]
\item No snake is warm-blooded. All mammals are warm-blooded. Therefore,
snakes are not mammals.

\begin{enumerate}
\item []\vspace{-25pt}%
\begin{tabular}{rll}
 &  & \tabularnewline
 &  & \tabularnewline
\hspace*{10pt} & \hspace*{2.5cm} & 1. \tabularnewline
\hspace*{10pt} & \hspace*{2.5cm} & 2.\hspace*{4.5cm}\tabularnewline
\cline{3-3}
 &  & \tabularnewline
\end{tabular}
\item [] \hspace{-1.25cm}%
\begin{tabular}{c|c|c|c|c|c|c|c|c|c}
\multicolumn{1}{c}{} & \multicolumn{1}{c}{} & \multicolumn{1}{c}{} & \multicolumn{1}{c}{} & \multicolumn{1}{c}{} & \multicolumn{1}{c}{} & \multicolumn{1}{c}{} & \multicolumn{1}{c}{} & \multicolumn{1}{c}{} & \tabularnewline
\multicolumn{1}{c}{} & \multicolumn{1}{c}{} & \multicolumn{1}{c}{} & \multicolumn{1}{c}{} & \multicolumn{1}{c}{} & \multicolumn{1}{c}{\ovalbox{$1$}} & \multicolumn{1}{c}{\ovalbox{$2$}} & \multicolumn{1}{c}{} & \multicolumn{1}{c}{\ovalbox{$c$}} & \ovalbox{Argument}\tabularnewline
 &  &  & \xspace & \xspace & \xspace & \xspace & \xspace & \xspace & \xspace\tabularnewline
\hline 
\hline 
\rowcolor{\rowColor}T & T & T &  &  &  &  &  &  & \tabularnewline
\hline 
T & T & F &  &  &  &  &  &  & \tabularnewline
\hline 
\rowcolor{\rowColor}T & F & T &  &  &  &  &  &  & \tabularnewline
\hline 
T & F & F &  &  &  &  &  &  & \tabularnewline
\hline 
\rowcolor{\rowColor}F & T & T &  &  &  &  &  &  & \tabularnewline
\hline 
F & T & F &  &  &  &  &  &  & \tabularnewline
\hline 
\rowcolor{\rowColor}F & F & T &  &  &  &  &  &  & \tabularnewline
\hline 
F & F & F &  &  &  &  &  &  & \tabularnewline
\end{tabular}\vfill
\end{enumerate}
\end{enumerate}
\end{enumerate}

\end{document}
